% STATUS: DRAFT
\chapter{Topology crash course}\label{app:topology}
\textbf{Open and closed sets.} A topology is a collection of ``open'' sets closed under arbitrary unions and finite intersections; closed sets are complements. A set is \emph{closed} iff it contains the limits of all convergent nets of its elements. \textbf{Nets.} In non-metrizable spaces (weak topologies on $\BH$) sequences are replaced by nets $(a_\alpha)$ indexed by a directed set. \textbf{Compactness.} $K$ is compact if every open cover has a finite subcover; equivalently every net has a convergent subnet. Compact subsets of Hausdorff spaces are closed. \textbf{Continuity.} $f$ is continuous iff preimages of open sets are open. \textbf{Weak topologies.} The weak topology induced by a family of functionals is the coarsest making them continuous; it is used for operators (\cref{ch:topologies}) and for states (weak-$^*$, \cref{ch:states}). \textbf{Banach--Alaoglu.} The closed unit ball of the dual of a Banach space is weak-$^*$ compact. \textbf{Baire/Hahn--Banach} are used implicitly: Hahn--Banach produces states and the singular states of \cref{ch:states}.
